\documentclass[10pt,a4paper]{amsart}
\usepackage[margin=19mm]{geometry}
\usepackage{amsmath,amssymb,mathtools}
\usepackage[scr=rsfs,cal=euler]{mathalfa}
\usepackage{xfrac}
\usepackage[unicode,hidelinks,bookmarks=false]{hyperref}
\newcommand{\ie}{i.e.\ }
\newcommand{\cf}{cf.\ }
\newcommand{\resp}{respectively\ }
\newcommand{\Subsection}{\subsection}
\providecommand{\nobreakdash}{\nobreak\hbox{-}}
\setlength{\emergencystretch}{2em}
\multlinegap0pt
\usepackage{bm}
\usepackage{bbm}
\usepackage{dsfont}
\usepackage{xspace}
\usepackage{cancel}
\usepackage{mathtools}
\usepackage{amsthm}
\makeatletter
\@ifundefined{theorem}{\newtheorem{theorem}{Theorem}}{}
\@ifundefined{lemma}{\newtheorem{lemma}[theorem]{Lemma}}{}
\makeatother
\makeatletter
\newcommand{\R}{\mathbb{R}}
\newcommand{\dd}{\,\mathrm{d}}
\renewcommand{\epsilon}{\varepsilon}
\expandafter\ifx\csname example\endcsname\relax
\newenvironment{example}
  {\pushQED{\qed}\renewcommand{\qedsymbol}{$\triangle$}\examplex}
  {\popQED\endexamplex}
\fi
\expandafter\ifx\csname remark\endcsname\relax
\newenvironment{remark}
  {\pushQED{\qed}\renewcommand{\qedsymbol}{$\triangle$}\remarkx}
  {\popQED\endremarkx}
\fi
\newcommand{\sequence}[2]{(#1_#2)_{#2}}
\newcommand{\veps}{\varepsilon}
\makeatother
\title[Homogenization of nonlocal exchange energies in micromagneti]{Homogenization of nonlocal exchange energies in micromagnetics}
\author{Rossella Giorgio, Leon Happ, Hidde Sch\"onberger}
\date{Source: arXiv:2507.13262v1 (CC BY 4.0)}
\begin{document}
\maketitle

\noindent\emph{Source and licence.} Homogenization of nonlocal exchange energies in micromagnetics. Version arXiv:2507.13262v1 (CC BY 4.0), distributed under the Creative Commons Attribution 4.0 International licence, as recorded in the arXiv licence metadata for this exact version and shown on the abstract page https://arxiv.org/abs/2507.13262v1. Full rights notice: licenses/arxiv-2507.13262-CC-BY-4.0.txt.\par\smallskip

\noindent\textit{Editorial scope.} Counted unit: the Lemma whose source label is \texttt{lem:x\_ts\_approx} in the article (source0.tex lines 877--884, source label \texttt{lem:x\_ts\_approx}), delivered together with the article's own complete original proof of it (source0.tex lines 885--922, one \texttt{proof} environment of 38 lines). No mathematics is written, repaired or re-derived: statement and proof are the article's own text line for line. Required context retained uncounted: lem:xts\_comp (lines 844--846). Every definition, display and label of those lines is retained unchanged. Omitted from this package and not redistributed: 1--843, 847--876, 923--2319. Nothing inside a retained excerpt is omitted or altered: every retained body is delivered exactly as the article's lines are written, so no omission receipt of any kind accompanies this package. Citations: the retained excerpts carry no citation command at all, so no bibliography entry of the article is transcribed and no reference is redistributed. Reference adaptation: none was needed and none was made; every reference inside the delivered unit resolves inside it, so no unresolved reference or citation is produced. Printed slips kept literal: no printed word, symbol, hypothesis or number of the counted statement or of its proof was altered. Layout and class: the article's own class is \texttt{amsart} with packages including geometry, amsmath, amssymb, amsthm, amscd, mathtools, url, fontenc, color, graphicx, enumitem, esint, dsfont, subfigure; the wrapper is the lane's pinned \texttt{amsart} editorial wrapper at 10pt on A4, which restates the theorem-like environments the retained text needs and adds the article's own macro definitions that the retained text needs (\texttt{\textbackslash R}, \texttt{\textbackslash dd}, \texttt{\textbackslash epsilon}, \texttt{\textbackslash sequence}, \texttt{\textbackslash veps}), with extra wrapper packages (bm, bbm, dsfont, xspace, cancel, mathtools). The delivered numbering is the unit's own. Assets: no retained span contains a figure environment, a graphics-inclusion command, a PDF-page inclusion, a table or a TikZ picture; no image file is copied into this package, the package declares no asset dependency and no image file is redistributed with it. Licence: version arXiv:2507.13262v1 of this article is distributed under the Creative Commons Attribution 4.0 International licence, as recorded in the arXiv licence metadata for this exact version and shown on the abstract page https://arxiv.org/abs/2507.13262v1. A full rights notice is delivered at licenses/arxiv-2507.13262-CC-BY-4.0.txt. Characters: every retained excerpt is pure ASCII, so the delivered engine, XeLaTeX with its default Latin Modern text font, reports no missing character.
\begin{lemma}\label{lem:xts_comp}
    For each bounded sequence \(\sequence{v}{\veps}\subset L^2(E;\R^m)\) there exists a subsequence and some \(v\in L^2(E_\xi\times E_x\times Q_x,\R^m)\) such that the subsequence weakly \(x\)-two-scale converges to \(v\).
\end{lemma}

\begin{lemma}\label{lem:x_ts_approx}
    Let \(\sequence{v}{\veps}\subset L^2(E;\R^m)\) be a bounded sequence and \(v\in L^2(E_\xi\times E_x\times Q_x;\R^m)\). Let \(D\subset L^2(E;C_{\#}(Q_x;\R^m))\) be a dense subset of functions in the respective topology. If 
    \begin{equation}\label{eq:x_ts_approx}
    \lim_{\veps \to 0}\int_{E_x}\int_{E_\xi} v_\epsilon(x,\xi) \cdot \psi(x,\xi,x/\epsilon)\dd \xi\dd x 
    = \int_{E_x}\int_{E_\xi}\int_{Q_x}v(x,\xi,z) \cdot \psi(x,\xi,z)\dd z\dd \xi \dd x
    \end{equation}
    for all \(\psi\in D\), then \(\sequence{v}{\veps}\) weakly \(x\)-two-scale converges to \(v\).
\end{lemma}

\begin{proof}
By Lemma~\ref{lem:xts_comp} we know that, up to a non-relabeled subsequence, $(v_\epsilon)_\epsilon$ weakly $x$-two-scale converges to some $v'\in L^2(E_\xi\times E_x\times Q_x;\R^m)$. Then, by \eqref{eq:x_ts_approx}, we deduce that
\[
\int_{E_x}\int_{E_\xi}\int_{Q_x}v(x,\xi,z) \cdot \psi(x,\xi,z)\dd z\dd \xi \dd x = \int_{E_x}\int_{E_\xi}\int_{Q_x}v'(x,\xi,z)\cdot \psi(x,\xi,z)\dd z\dd \xi \dd x,
\]
for all $\psi \in D$. By density of $D$, we infer that $v=v'$. Hence, up to a subsequence, $(v_\epsilon)_\epsilon$ weakly $x$-two-scale converges to $v$. Since this limit is independent of the chosen subsequence, the convergence holds for the entire sequence.
    % We adapt the proof strategy of \cite[Proposition~1]{LNW02}. For any \(\psi \in L^2(E;C_{\#}(Q_x;\R^3m))\) and a sequence \(\seq{\psi}\subset D\) converging to \(\psi\) in \(L^2(E;C_{\#}(Q_x;\R^m))\) there holds
    % \begin{align*}
    %     &\iint_{E} v_\epsilon(x,\xi) \psi(x,\xi,x/\epsilon)\dd x\dd \xi \\
    %     &\quad = \iint_{E} v_\epsilon(x,\xi) \psi_k(x,\xi,x/\epsilon)\dd x\dd \xi 
    %     +\iint_{E} v_\epsilon(x,\xi) \left(\psi(x,\xi,x/\epsilon)-\psi_k(x,\xi,x/\epsilon)\right)\dd x\dd \xi .
    % \end{align*}
    % The first integral tends to 
    % \begin{equation}\label{eq:x_ts_approx_aux1}
    %     \iiint_{E}v(x,\xi,z)\psi_k(x,\xi,z)\dd z\dd x \dd \xi
    % \end{equation}
    % for \(\veps\to 0\), following assumption \eqref{eq:x_ts_approx}, whilst the second one might be estimated with the help of H{\"o}lder's inequality and the boundedness of \(\sequence{v}{\veps}\) by
    % \begin{align}
    % \begin{split}\label{eq:x_ts_approx_aux2}
    %     &\Bigg|\iint_{E} v_\epsilon(x,\xi) \left(\psi(x,\xi,x/\epsilon)-\psi_k(x,\xi,x/\epsilon)\right)\dd x\dd \xi \Bigg|\\
    %     &\quad\le \norm{v_\veps}_{L^2(E;\R^m)}
    %     \left(\iint_{E} \Big|\psi(x,\xi,x/\epsilon)-\psi_k(x,\xi,x/\epsilon)\Big|^2\dd x\dd \xi\right)^{1/2}\\
    %     &\quad\le C\norm{\psi-\psi_k}_{L^2(E;C_{\#}(Q_x;\R^m))}.
    % \end{split}
    % \end{align}
    % Letting now \(k\to \infty\), we find on the one hand in view of \eqref{eq:x_ts_approx_aux1} that
    % \begin{equation*}
    %     \lim_{k\to\infty}\iiint_{E_\xi\times E_x\times Q_x}v(x,\xi,z)\psi_k(x,\xi,z)\dd z\dd x \dd \xi
    %     = \iiint_{E_\xi\times E_x\times Q_x}v(x,\xi,z)\psi(x,\xi,z)\dd z\dd x \dd \xi,
    % \end{equation*}
    % since
    % \begin{align*}
    %     &\Bigg|\iiint_{E_\xi\times E_x\times Q_x}v(x,\xi,z)(\psi(x,\xi,z)-\psi_k(x,\xi,z))\dd z\dd x \dd \xi\Bigg|\\
    %     &\quad\le \norm{v}_{L^2(E_\xi\times E_x\times Q_x;\R^m)}\norm{\psi-\psi_k}_{L^2(E_\xi\times E_x\times Q_x;\R^m)}\\
    %     &\quad\le \norm{v}_{L^2(E_\xi\times E_x\times Q_x;\R^m)}\norm{\psi-\psi_k}_{L^2(E;C_{\#}(Q_x;\R^m))}
    % \end{align*}
    % and \(\seq{\psi}\) converges to \(\psi\) in \(L^2(E;C_{\#}(Q_x;\R^m))\). On the other hand, by reasoning similarly, the remainder in \eqref{eq:x_ts_approx_aux2} vanishes for \(k\to \infty\). This concludes the proof.
\end{proof}

\end{document}
